\section{Relativistic layers: scalar X2C, two-component X2C, and ECP spin--orbit}
\label{sec:relativistic-tutorial}

This chapter runs the three relativistic entry points that do not
belong to the differentiable multireference chain: the scalar
(spin-free) density-fitted X2C CASSCF, the two-component density-fitted
X2C Hartree--Fock, and the valence-level spin--orbit mean field for
effective core potentials.  All three are NumPy or pyscf layers; none
runs on pyscfad and none contributes to the JAX derivative chain.

\subsection{Scalar DF-X2C CASSCF}

The scalar layer transforms the one-electron Hamiltonian with the
spin-free Foldy--Wouthuysen operator and keeps the non-relativistic
Coulomb interaction, density fitted.  The wavefunction is a real
spin-adapted determinant expansion, so the CASSCF machinery is the same
as in the non-relativistic chapters.

\begin{lstlisting}[language=Python]
import numpy as np
from pyscf import gto
from diffmrc.core import x2c_casscf_df as xc

mol = gto.M(atom='O 0 0 0; H 0 -0.757 0.587; H 0 0.757 0.587',
            basis='6-31g', unit='Angstrom', verbose=0)

res = xc.df_x2c_casscf(mol, ncore=4, nact=2, nele=2,
                       auxbasis='def2-universal-jkfit',
                       g_conv=1e-6, max_cycle=400)
print('E', res.e_tot, 'gnorm', res.gnorm,
      'converged', res.converged)
print('aux basis', res.auxbasis)

grad, parts = xc.df_x2c_casscf_gradient(mol, res)
print('gradient', grad)
for name, value in parts.items():
    print(f'{name:14s} {np.abs(value).max(): .3e}')
\end{lstlisting}

The four \texttt{parts} are the one-electron X2C term, the
density-fitted two-electron term, the Pulay term, and the nuclear
repulsion.  Their sum is the gradient; the split is useful when a
gradient looks wrong, because a factor-derivative error and a Pulay
sign error do not have the same size.

The single-state gradient is the direct derivative at fixed orbitals
and CI vector: at a converged CASSCF point both responses vanish.  A
state-averaged calculation is different, because the orbitals minimize
the average and the target state is not stationary:

\begin{lstlisting}[language=Python]
res_sa = xc.df_x2c_casscf(mol, ncore=4, nact=2, nele=2,
                          auxbasis='def2-universal-jkfit',
                          nroots=2, weights=[0.5, 0.5],
                          g_conv=1e-6, max_cycle=400)
print('SA energy', res_sa.e_sa, 'states', res_sa.state_energies)

g0, info = xc.df_x2c_sa_state_gradient(mol, res_sa, state=0,
                                       return_info=True)
print('state 0 gradient', g0)
print('response directions', info['n_response'])
print('retained directions', info['n_kept'])
print('Hessian spectrum', info['hess_eigvals'])
\end{lstlisting}

The state-averaged response is a finite-difference construction of the
orbital Hessian and the mixed derivative around the analytic
state-averaged orbital gradient.  The CI vectors are re-diagonalized at
every displacement, so the CI response is inside the Hessian.  This is
the route to use when a per-state gradient is needed; do not replace it
by re-optimizing the orbitals at the displaced geometry, which on water
can land in a second orbital minimum.

\subsection{Two-component DF-X2C Hartree--Fock}

The two-component layer keeps the spinor structure of the X2C
transform.  The Hamiltonian is complex Hermitian, the SCF is a complex
generalized Hartree--Fock, and the density-fitted Coulomb and exchange
matrices are built directly in the contracted spinor basis.

\begin{lstlisting}[language=Python]
import numpy as np
from pyscf import gto
from pyscf.x2c import x2c as pyscf_x2c
from diffmrc.core import x2c_df

mol = gto.M(atom='O 0 0 0; H 0 -0.757 0.587; H 0 0.757 0.587',
            basis='6-31g', unit='Angstrom', verbose=0)

h = x2c_df.df_x2c_hcore(mol)
h_ref = pyscf_x2c.SpinorX2CHelper(mol).get_hcore(mol)
print('hcore max difference', np.abs(h - h_ref).max())

hf = x2c_df.df_x2c_hf(mol, auxbasis='def2-universal-jkfit')
print('E', hf.e_tot, 'gnorm', hf.gradient_norm,
      'converged', hf.converged)
print('orbital energies', hf.mo_energy[:4])

grad, parts = x2c_df.df_x2c_hf_gradient(mol, result=hf)
print('gradient', grad)
print('parts', {k: float(np.abs(v).max()) for k, v in parts.items()})
\end{lstlisting}

The one-electron check against pyscf is the first thing to run on a
new basis: it isolates the X2C transform from the density fitting.  If
the Hamiltonian differs, the SCF and the gradient will both be wrong
for a reason that has nothing to do with the two-electron part.

The spinor ERI convention is the complex-GHF one.  The three-center
tensor is Hermitian in the AO pair, and the derivative with respect to
the second AO center carries a complex conjugate.  The test suite pins
that conjugate against a geometry finite difference; a gradient that
looks right at the $10^{-3}$ level but not at $10^{-6}$ usually has a
missing conjugate or a wrong AO slice.

\subsection{Valence-level SOMF for ECP calculations}

The all-electron SOMF operator uses the bare nuclear charges and the
full electron density.  With an effective core potential the core
electrons are absent, and the spin--orbit part of the ECP replaces
their contribution.  The valence-level operator is
\begin{equation}
  h^{\mathrm{SOMF\text{-}ECP}}_k
  = \frac{\alpha^{2}}{2}\left[h^{\mathrm{nuc}}_k + h^{2e}_k[P]\right]
    - V^{\mathrm{ECP}}_k,
  \label{eq:tutorial-somf-ecp}
\end{equation}
with the nuclear term kept only on centers without a spin--orbit ECP,
the valence $J-\tfrac{3}{2}K$ mean field of the spin-summed density
$P$, and the three-component \texttt{ECPso} integral.

\begin{lstlisting}[language=Python]
import numpy as np
from pyscf import gto
from diffmrc.core import somf_ecp

mol = gto.M(atom=[('Ge', (0.0, 0.0, 0.0)), ('F', (0.5, 0.3, 2.4))],
            basis={'Ge': 'lanl2dz', 'F': '6-31g'},
            ecp={'Ge': 'ecpds28mdfso'}, spin=1,
            unit='Bohr', verbose=0)

# A density of the electrons present in the calculation.  Here it is the
# equal-occupation reference density; a converged SCF density works too.
s = mol.intor('int1e_ovlp')
P = (mol.nelectron / mol.nao_nr()) * np.linalg.inv(s)
P = 0.5 * (P + P.T)

h_soc = somf_ecp.valence_somf_matrix(mol, P, convention='i', check=True)
print('SOC matrix shape', h_soc.shape)
print('antisymmetry', np.abs(h_soc + h_soc.transpose(0, 2, 1)).max())

# The density response is a separate contraction; the SOC kernel is not
# symmetric under (pq) <-> (rs).
m_resp = somf_ecp.somf_density_response(mol, P)
print('response shape', m_resp.shape)

print(somf_ecp.somf_definition(mol))
\end{lstlisting}

\texttt{somf\_definition} records which centers contribute the ECP
term, which keep the nuclear term, and the approximations that the ECP
data do not fix.  The operator is a prototype for relative comparisons,
not a production absolute spin--orbit method: the split between the
ECP spin--orbit operator and the explicit valence mean field is a
fitting freedom.  On the GeF model with a scalar ECP the valence-level splitting is
$26.00$~cm$^{-1}$ against $40.14$~cm$^{-1}$ for the nuclear-only
operator; report the variant bracket rather than a single number.

The eligibility rule is simple: the molecule must carry at least one
spin--orbit ECP (\texttt{mol.has\_ecp\_soc()} is true).  A scalar-only
ECP does not supply \texttt{ECPso}, and \texttt{ecp\_soc} raises.  On
centers without a spin--orbit ECP the nuclear term is used with the
effective charge from the atom table.

\subsection{Where the scalar and two-component X2C energies differ}

The scalar layer and the two-component layer are not a controlled
one-parameter pair: the scalar layer is CASSCF and the two-component
layer is HF.  At the same wavefunction level the X2C-specific
difference separates into three terms,
\begin{equation}
  E_{2c}-E_{1c}
  = \Delta E_{\mathrm{spin\text{-}free}}
  + \Delta E_{\mathrm{SO}}
  + \Delta E_{2e},
  \label{eq:tutorial-x2c-difference}
\end{equation}
where $\Delta E_{\mathrm{spin\text{-}free}}$ is the change from the
spin-dependent renormalization and picture-change matrix that the
spin-free transform drops, $\Delta E_{\mathrm{SO}}$ is the
spin--orbit contribution, and $\Delta E_{2e}$ is the difference
between the spatial DF and spinor RI-JK two-electron treatments.

For a non-degenerate closed-shell determinant the first-order
spin--orbit expectation value is quenched:
\begin{equation}
  \langle \Phi | H_{\mathrm{SO}} | \Phi \rangle = 0
  \qquad \text{(closed-shell, non-degenerate)}.
  \label{eq:tutorial-soc-quenched}
\end{equation}
The surviving $\Delta E_{\mathrm{SO}}$ is second order and comes from
the response of the orbitals and the CI, not from a direct matrix
element.  For an open-shell or degenerate system the first-order term
is not quenched and can dominate.  The measured OH/OF splittings in the
manual chapter are an open-shell example; a closed-shell water
comparison is dominated by the spin-free and two-electron terms.

The practical rule is: do not quote a scalar-versus-two-component
energy difference as a relativistic correction unless the wavefunction
method and the two-electron operator are the same on both sides.  The
shipped scalar CASSCF and two-component HF do not satisfy that
condition.

\subsection{Which layer to use}

\begin{center}
\small
\begin{tabularx}{\textwidth}{>{\raggedright\arraybackslash}p{0.25\textwidth}X}
\toprule
Situation & Layer \\
\midrule
Real, spin-free multireference calculation on a full-electron basis &
  \texttt{core.x2c\_casscf\_df} \\
Complex one-electron spin--orbit at the HF level &
  \texttt{core.x2c\_df} \\
Spin--orbit splitting of a valence configuration with an ECP &
  \texttt{core.somf\_ecp} (relative comparisons) \\
Two-electron picture-change correction &
  \texttt{core.x2c\_sap} (separate, not covered here) \\
ECP multireference with spin--orbit coupling &
  not in this release \\
\bottomrule
\end{tabularx}
\end{center}
